\documentclass[12pt]{article}

\usepackage[a4paper, margin=2.5cm]{geometry}
\usepackage{amsmath}
\usepackage{amssymb}
\usepackage{amsthm}
\usepackage[colorlinks=true, linkcolor=blue, citecolor=blue, urlcolor=blue]{hyperref}
\usepackage{booktabs}
\usepackage{natbib}
\usepackage{url}

\newtheorem{theorem}{Theorem}
\newtheorem{proposition}[theorem]{Proposition}
\newtheorem{corollary}[theorem]{Corollary}
\newtheorem{definition}[theorem]{Definition}
\newtheorem{remark}[theorem]{Remark}
\newtheorem{lemma}[theorem]{Lemma}

\newcommand{\phig}{\varphi}
\newcommand{\lbar}{\bar{\lambda}}
\newcommand{\Fcal}{\mathcal{F}}
\newcommand{\Mcl}{\mathcal{M}_{\mathrm{cl}}}
\newcommand{\Sadm}{\mathcal{S}_{\mathrm{adm}}}
\newcommand{\Scl}{\mathcal{S}_{\mathrm{cl}}}
\newcommand{\HW}{H_{\mathrm{W}}}
\newcommand{\VW}{V_{\mathrm{W}}}

\title{Hadron Charge Radii as Spectral Coherence Lengths\\[6pt]
\large from 600-Cell Closure Geometry}
\author{%
  Lee Smart\\[4pt]
  \textit{Institute of Vibrational Field Dynamics}\\[2pt]
  \texttt{contact@vibrationalfielddynamics.org}\\[2pt]
  \texttt{@vfd\textunderscore org}%
}
\date{April 2026}

\begin{document}

\maketitle

\noindent\textbf{Status:} Preprint (not peer-reviewed)

\begin{abstract}
We derive the charge radii of six hadrons from the $\phig$-structured 600-cell closure geometry, with zero fitted parameters. The central result is a unified principle: every charge radius is the heat-kernel coherence length of the particle's electromagnetic response, computed on the geometric sub-object of the 600-cell where the particle lives. The coherence mechanism depends on the particle's structural class: baryons use the boundary resolvent on their shell support graph, non-strange mesons use the phase coherence of the closure orbit, and strange mesons use the diffusion length on the Hopf fiber.

The proton charge radius is derived from the boundary resolvent on the proton support graph $P_3$. The resulting coherence length can be written as $r_p = 4\,\hbar/(m_p c)$, where the factor $4$ equals both the Laplacian trace $\mathrm{Tr}(L(P_3))$ and $\max(S_p)$; the latter equality is proved as a theorem of the assignment rules. The result $r_p = 0.8412$~fm agrees with the PDG value $0.8409 \pm 0.0004$~fm at 0.04\%. Five additional radii are derived: neutron $\langle r_n^2 \rangle = -(8/3)\,\lbar_n^2$ (1.3\%), pion $r_\pi = \pi\,\lbar_p$ (0.26\%), kaon $r_K = \phig^2\,\lbar_p$ (1.7\%), proton magnetic $r_M = r_E$ (prediction), and deuteron $r_d \approx 2.126$~fm from a three-component geometric decomposition (0.10\% relative error; outside $1\sigma$ because the experimental uncertainty is very small). The kaon result is derived from the spectral gap of the Hopf decagonal fiber $C_{10}$: $E_1(C_{10}) = \phig^{-2}$, connecting the charge radius to the generation structure.

The proton electric form factor is derived as the rational function $G_E(Q^2) = (1+3z+z^2)/(1+4z+3z^2)$, with timelike zeros at the golden-ratio scales $z = -\phig^{-2}$ and $z = -\phig^2$. The 120-vertex resolvent form factor matches low-$Q^2$ data within 0.4\% at $Q^2 = 0.1$~GeV$^2$, with deviations growing at higher $Q^2$.

No continuous parameters are fitted. Every quantity is either a fundamental constant ($\hbar$, $c$, $m_e$) or a structural invariant of the 600-cell geometry ($\phig$, eigenvalues, graph invariants). Six observables are derived within the framework from one geometric principle, with zero free parameters.
\end{abstract}

%==================================================================
\section{Introduction}\label{sec:intro}
%==================================================================

The charge radii of hadrons are fundamental observables encoding the spatial distribution of electromagnetic charge. The proton charge radius, defined via the slope of the electric form factor:
\begin{equation}\label{eq:radius_def}
    r_p^2 = -6\,\frac{dG_E(q^2)}{dq^2}\bigg|_{q^2=0}
\end{equation}
is measured to high precision~\citep{Workman2022}: $r_p = 0.8409 \pm 0.0004$~fm. In the Standard Model, this radius is computed from lattice QCD~\citep{Durr2008} or extracted from electron--proton scattering and muonic hydrogen spectroscopy. The question of whether hadron charge radii can be derived from geometric or algebraic structure, rather than fitted to experiment, remains open.

The present paper derives six hadron charge radii from the VFD closure framework~\citep{SmartIV, SmartVIII, SmartXXII}, which constructs particle masses, coupling constants, and gauge structure from the spectral geometry of the 600-cell polytope. The results use zero fitted continuous parameters and agree with experiment at the sub-2\% level for all six observables.

The derivation rests on a single principle: every charge radius is the \textit{heat-kernel coherence length} of the particle's electromagnetic response, computed on the geometric sub-object of the 600-cell where the particle lives. The coherence mechanism depends on the particle's structural class:
\begin{itemize}
    \item \textbf{Baryons} are closure basins on the shell support graph $G(S)$. Their radii are determined by the boundary resolvent on the support graph.
    \item \textbf{Non-strange mesons} are phase excitations along the closure orbit. Their radii are determined by the orbital phase coherence length.
    \item \textbf{Strange mesons} are cross-generation modes on the Hopf fibers. Their radii are determined by the fiber spectral gap.
\end{itemize}

All three geometric objects --- the support graph $P_3$, the closure orbit $S^1$, and the Hopf fiber $C_{10}$ --- live inside the same 600-cell. The heat kernel on the 600-cell is the master object from which all coherence lengths are extracted.

\subsection*{Scope and epistemic status}

This paper derives charge radii, not internal structure. The results are structural: they follow deterministically from the 600-cell geometry and the closure assignment rules, without optimisation or parameter adjustment. The paper does not claim to replace QCD, derive the Standard Model from first principles, or explain quark confinement dynamically. It shows that, within the VFD closure framework, six independent charge radii emerge from one geometric principle with sub-2\% accuracy. The honest classification of each result's derivation status is given throughout.

%==================================================================
\section{Framework Essentials}\label{sec:framework}
%==================================================================

We summarise the minimal framework elements used in the derivation. Full details are in Papers~I--XXII~\citep{SmartI, SmartIV, SmartVIII, SmartXXII}.

\subsection{The 600-cell and the golden ratio}

The 600-cell is a regular polytope in $\mathbb{R}^4$ with 120 vertices, 720 edges, and vertex degree 12. Its vertices are the elements of the binary icosahedral group $2I \subset \mathrm{SU}(2)$~\citep{Coxeter1973}. The golden ratio $\phig = (1+\sqrt{5})/2$ governs the geometry: vertex coordinates involve $\phig$, the nearest-neighbour distance is $1/\phig$, and the spectral gap of the shell path graph $P_5$ is $\lambda_1(P_5) = 2 - 2\cos(\pi/5) = \phig^{-2}$.

\subsection{Shell support and assignment rules}

A $\phig$-scaled shell manifold $\mathcal{M}$ consists of nested shells at scales $r_n = r_0\,\phig^{-n}$ for $n = 1, \ldots, 5$. A particle occupies a shell support $S \subseteq \{1, \ldots, 5\}$ with support graph $G(S)$: one vertex per occupied shell, one edge between consecutive shells.

Five rules~\citep{SmartIV} constrain admissible supports:
\begin{itemize}
    \item[\textbf{R1}] Minimal support: composites require $|S| \geq 3$.
    \item[\textbf{R2}] Boundary/interior: baryons start at shell 2.
    \item[\textbf{R3}] Complexity threshold: $C(S) \geq 5$ for baryons.
    \item[\textbf{R4}] Contiguity: shells must be consecutive.
    \item[\textbf{R5}] Connectivity: $G(S)$ must be connected.
\end{itemize}
Under R1--R5, the electron is uniquely assigned to $S_e = \{1\}$ and the proton to $S_p = \{2,3,4\}$~\citep{SmartIV}.

\subsection{The proton support graph}

The proton support graph $G(S_p) = P_3$ is a path on 3 vertices with:
\begin{itemize}
    \item $|V| = 3$ vertices, $|E| = 2$ edges
    \item Degree sequence $[1, 2, 1]$
    \item Laplacian eigenvalues $\{0, 1, 3\}$
    \item Trace $\mathrm{Tr}(L) = 0 + 1 + 3 = 4$
    \item Degree variance $\mathrm{Var}(\deg) = 2/9$
\end{itemize}

\subsection{The Hopf fiber decomposition}

The 120 vertices of the 600-cell decompose under the Hopf fibration $S^3 \to S^2$ into 12 decagonal great circles $C_{10}$, each passing through 10 vertices. The eigenvalues of $C_{10}$ are $E_k = 2 - 2\cos(2\pi k/10)$; the spectral gap is:
\begin{equation}\label{eq:fiber_gap}
    E_1(C_{10}) = 2 - 2\cos(\pi/5) = \phig^{-2}
\end{equation}
Generations correspond to winding modes on the Hopf fiber~\citep{SmartXXII}: $w = 1$ (first generation), $w = 2$ (second), $w = 3$ (third).

\subsection{The Compton wavelength as step length}

In the tight-binding regime~\citep{SmartXXII}, the Witten Hamiltonian reduces to $\HW^{\mathrm{eff}} \approx E_0\,\mathbf{I} + t\,\mathbf{A}_{600}$ with uniform tunnelling by vertex-transitivity. Each tunnelling step contributes one reduced Compton wavelength $\lbar = \hbar/(mc)$ to the spatial extent (standard WKB result at the mass scale $m$). The proton's Compton wavelength $\lbar_p = \hbar/(m_p c) = 0.210309$~fm is the natural length unit.

%==================================================================
\section{The Heat Kernel from the Closure Functional}\label{sec:heatkernel}
%==================================================================

Within the tight-binding reduction established in the programme, the relevant heat kernel is induced by the closure functional through the effective Laplacian. The chain of constructions is explicit.

\subsection{From $\Fcal$ to the Laplacian}

The closure functional on the 600-cell~\citep{SmartXXII}:
\begin{equation}\label{eq:closure_functional}
    \Fcal(x) = -\varepsilon^2 \log\!\left(\sum_{i=1}^{120} \exp\!\left(-\frac{|x - v_i|^2}{2\varepsilon^2}\right)\right)
\end{equation}
has uniform Hessian $H_\Fcal(v_j) = \varepsilon^{-2}\,\mathbf{I}$ at each vertex (by vertex-transitivity). The Witten potential $\VW = |\nabla\Fcal|^2/(2\sigma^2) - \Delta\Fcal/2$ is constant on the vertex set for the same reason. In the tight-binding regime ($\varepsilon \ll d_1$), the Witten Hamiltonian reduces to~\citep{SmartXXII}:
\begin{equation}
    \HW^{\mathrm{eff}} = E_0\,\mathbf{I} + t\,\mathbf{A}_{600}
\end{equation}
where $\mathbf{A}_{600}$ is the adjacency matrix and $t \propto \exp(-d_1^2/(2\sigma^2\varepsilon^2))$ is the uniform tunnelling amplitude. The graph Laplacian is $L = 12\,\mathbf{I} - \mathbf{A}_{600}$, so $\HW^{\mathrm{eff}} = (E_0 + 12t)\,\mathbf{I} - t\,L$.

Up to the constant shift, the effective Witten Hamiltonian is proportional to the graph Laplacian: $\HW^{\mathrm{eff}} \propto -L$.

\subsection{From the Laplacian to the heat kernel}

The heat kernel on the graph is:
\begin{equation}\label{eq:heatkernel}
    K(v_i, v_j; \tau) = \bigl(e^{-\tau L}\bigr)_{ij} = \sum_k e^{-\lambda_k \tau}\, \psi_k(v_i)\,\psi_k(v_j)
\end{equation}
where $\{\lambda_k, \psi_k\}$ are the eigenpairs of $L$. This is a standard construction: $K$ is the matrix exponential of $-\tau L$, which is the solution to $\partial_\tau K = -L\,K$ with $K(0) = \mathbf{I}$.

\subsection{The complete chain}

\begin{equation}
    \Fcal(x) \;\xrightarrow{\text{Hessian}}\; \HW \;\xrightarrow{\text{tight-binding}}\; L \;\xrightarrow{\text{matrix exponential}}\; K(v_i, v_j; \tau)
\end{equation}

Every step is either a standard construction (Hessian, matrix exponential) or an established result of the programme (tight-binding reduction, Paper~XXII). Within this reduction, the heat kernel is induced by $\Fcal$ through the effective Laplacian.

\subsection{Observables from the heat kernel}

The heat kernel encodes all spectral information:
\begin{itemize}
    \item \textbf{Masses}: the spectral gap $\lambda_1$ extracted from $\ln K \sim -\lambda_1 \tau$ at large $\tau$ (Paper~IV).
    \item \textbf{Coupling constants}: the eigenvalue algebra of $L$ (Paper~XXII).
    \item \textbf{Charge radii}: the coherence length $\xi_K$ (this paper).
    \item \textbf{Form factors}: the resolvent $R(z) = \int_0^\infty K(\tau)\,e^{-z\tau}\,d\tau$ (this paper).
\end{itemize}

The pipeline $\Fcal \to L \to K \to \text{observables}$ is closed and self-contained.

%==================================================================
\section{The Spectral Channel Operator}\label{sec:channel}
%==================================================================

The charge radius is extracted from the heat kernel through a \textit{coherence factor} $\xi_K$ that depends on the geometric home $G$ and a spectral channel $c$.

\begin{definition}[Spectral coherence radius]\label{def:channel}
For a particle class with geometric home $G$ and spectral channel $c$, let $\xi_K(G, c)$ denote the \textit{dimensionless coherence factor} extracted from the heat kernel on $G$ in channel $c$. The charge radius is:
\begin{equation}\label{eq:channel_r}
    r = \xi_K(G, c) \times \lbar_p \qquad\text{(mesonic channels)}
\end{equation}
\begin{equation}\label{eq:channel_r2}
    r^2 = \xi_K(G, c) \times \lbar_p^2 \qquad\text{(baryonic resolvent channels)}
\end{equation}
where $\lbar_p = \hbar/(m_p c)$ is the confinement Compton wavelength.
\end{definition}

The three channels extract different spectral functionals:

\begin{table}[ht]
\centering
\caption{Spectral channels and their dimensionless coherence factors.}
\label{tab:channels}
\begin{tabular}{@{}lllll@{}}
\toprule
Channel $c$ & $\xi_K$ & Spectral origin & Particles & Example \\
\midrule
Monopole resolvent & $12\,\mathrm{Tr}(L)/|V|$ & Total boundary response & Charged baryons & $\xi = 16$, $r = 4\,\lbar$ \\
Dipole resolvent & $\mathrm{Var}(\deg)\,\Delta n^2$ & Charge separation & Neutral baryons & $\xi = 8/3$ \\
Phase coherence & $\pi$ & Orbital half-period & Goldstone mesons & $r = \pi\,\lbar$ \\
Spectral diffusion & $\lambda_1^{-1}$ & Inverse spectral gap & Strange mesons & $r = \phig^2\,\lbar$ \\
\bottomrule
\end{tabular}
\end{table}

These are the three canonical spectral functionals of a graph Laplacian: the trace (monopole), the first eigenfunction (dipole), and the spectral gap (diffusion). The phase channel is the continuum limit of the diffusion channel on the closure orbit.

\begin{remark}[Channel selection is structural]
The channel is determined by the particle's charge structure, not chosen post hoc:
\begin{itemize}
    \item Total charge $Q \neq 0$ $\Rightarrow$ monopole resolvent (total boundary response).
    \item Total charge $Q = 0$ $\Rightarrow$ dipole resolvent (charge cancels in monopole; first nontrivial eigenfunction carries the residual).
    \item Pseudo-Goldstone boson $\Rightarrow$ phase or diffusion (correlation length sets the spatial extent).
\end{itemize}
\end{remark}

%==================================================================
\section{Proton Charge Radius}\label{sec:proton}
%==================================================================

\subsection{The boundary resolvent}

The baryon charge radius is extracted from the resolvent of the graph Laplacian $L(G(S))$ evaluated at the boundary vertex. For the proton support graph $P_3$, the resolvent at the outer boundary (shell~4, vertex~2) gives a charge radius:
\begin{equation}\label{eq:boundary}
    r_p^2 = 12\,\frac{\mathrm{Tr}(L)}{|V|}\,\lbar_p^2
\end{equation}
where $\mathrm{Tr}(L)/|V| = 2|E|/|V|$ is the mean degree --- the first nontrivial cumulant coefficient of the heat-kernel expansion (Paper~IV, Section~2.3). For $P_3$: $\mathrm{Tr}(L)/|V| = 4/3$, giving $r_p = \sqrt{16}\,\lbar_p = 4\,\lbar_p$.

\begin{remark}[$P_3$-specific simplification]
The numerical coincidence $\sqrt{12 \times 4/3} = 4 = \mathrm{Tr}(L(P_3))$ holds because the proton support satisfies $n_{\min} = |S| - 1$ (Proposition~\ref{prop:identity}). For general path graphs $P_n$, the resolvent formula~\eqref{eq:boundary} and $\mathrm{Tr}(L) \times \lbar$ are not equivalent; the resolvent formula is primary.
\end{remark}

\subsection{The structural identity}

\begin{proposition}\label{prop:identity}
$\mathrm{Tr}(L(P_3)) = \max(S_p)$ is structurally forced by R1--R5.
\end{proposition}

\begin{proof}
For a path graph $P_n$: $\mathrm{Tr}(L) = 2(n-1)$. For a contiguous support starting at $n_{\min}$ with $|S|$ shells: $\max(S) = n_{\min} + |S| - 1$. These are equal iff $n_{\min} = |S| - 1$. For the proton: R2 gives $n_{\min} = 2$, R1 gives $|S| \geq 3$, and minimality gives $|S| = 3$. Then $n_{\min} = 2 = 3 - 1 = |S| - 1$. Therefore $\mathrm{Tr}(L(P_3)) = 2(3-1) = 4 = 2 + 3 - 1 = \max(S_p)$.
\end{proof}

\subsection{Derivation from the boundary resolvent}\label{sec:resolvent}

The proton form factor is derived from the resolvent of the graph Laplacian at the boundary vertex (shell~4, vertex~2 of $P_3$). The eigenvector components at this vertex are $|\psi_k|^2 = \{1/3, 1/2, 1/6\}$ for $\lambda_k = \{0, 1, 3\}$. The resolvent:
\begin{equation}
    R_{\mathrm{bdy}}(z) = \frac{1/3}{1} + \frac{1/2}{1+z} + \frac{1/6}{1+3z}
\end{equation}
gives the form factor~\eqref{eq:formfactor} (Section~\ref{sec:formfactor}). The slope at $z = 0$: $dR/dz|_0 = -1$, yielding $r^2 = 6a^2$ with lattice spacing $a^2 = 2\,\mathrm{Tr}(L)/|V| \times \lbar^2 = 8\lbar^2/3$. Therefore:
\begin{equation}
    r_p^2 = 6 \times \frac{8\lbar_p^2}{3} = 16\,\lbar_p^2, \qquad r_p = 4\,\lbar_p = 0.8412~\text{fm}
\end{equation}

\subsection{Boundary dominance}\label{sec:boundary_dominance}

The electromagnetic probe couples to the \textit{boundary} of the closure configuration because interior vertices are shielded by the closure barrier. In the tight-binding regime, penetrating from the exterior to vertex $k$ requires traversing $k$ closure barriers, each suppressing the coupling by $\exp(-F_{\mathrm{barrier}}/\sigma^2)$ where $F_{\mathrm{barrier}} = d_1^2/(4\varepsilon^2)$. The outermost vertex (shell~4) has zero barriers; interior vertices are exponentially suppressed. This is confirmed by the eigenvector structure: the dipole mode $v_1 = [1, 0, -1]/\sqrt{2}$ has weight $1/2$ at each boundary vertex and exactly $0$ at the interior.

\subsection{Numerical result}

\begin{table}[ht]
\centering
\caption{Proton charge radius.}
\label{tab:proton}
\begin{tabular}{@{}lll@{}}
\toprule
Source & $r_p$ (fm) & Uncertainty \\
\midrule
\textbf{VFD (this work)} & \textbf{0.8412} & \textbf{structural} \\
PDG 2022 & 0.8409 & 0.0004~fm \\
Muonic hydrogen & 0.84087 & 0.00039~fm \\
\bottomrule
\end{tabular}
\end{table}

\noindent Relative error: 0.04\%. Within 0.84$\sigma$ of the PDG value. The proton magnetic radius is predicted to equal the charge radius: $r_M = r_E = 0.841$~fm (experimental: $0.851 \pm 0.026$~fm, within $0.4\sigma$).

%==================================================================
\section{Neutron Charge Radius}\label{sec:neutron}
%==================================================================

The neutron shares support $\{2,3,4\}$ with the proton. Its total charge is zero, so the monopole channel (used for the proton) gives zero. The charge radius comes from the \textit{dipole channel} --- the first nontrivial Laplacian eigenfunction of $P_3$.

\subsection{The dipole eigenfunction}

The eigenfunction $v_1 = [1, 0, -1]/\sqrt{2}$ (eigenvalue $\lambda_1 = 1$) places positive charge at shell~2 and negative charge at shell~4, with zero at the interior. The charge amplitude is:
\begin{equation}
    q = \mathrm{Var}(\deg) = \frac{2}{9}
\end{equation}
the degree variance of $P_3$. This is the same quantity that enters the second-order mass correction~\citep{SmartIV}.

\subsection{Result}

\begin{equation}\label{eq:neutron}
    \langle r_n^2 \rangle = -q \times (n_{\max}^2 - n_{\min}^2) \times \lbar_n^2 = -\frac{2}{9} \times 12 \times \lbar_n^2 = -\frac{8}{3}\,\lbar_n^2 = -0.1176~\text{fm}^2
\end{equation}

\noindent Experimental: $\langle r_n^2 \rangle = -0.1161 \pm 0.0022$~fm$^2$. Error: 1.3\%, within 0.69$\sigma$.

%==================================================================
\section{Pion Charge Radius}\label{sec:pion}
%==================================================================

The pion is a pseudo-Goldstone boson --- a phase excitation along the flat direction of the closure functional. Its charge radius measures the spatial extent of the phase coherence.

The closure orbit on $S^3$ has angular period $2\pi$. The pion's coherence length is the half-period:
\begin{equation}\label{eq:pion}
    r_\pi = \pi \times \lbar_p = 0.6607~\text{fm}
\end{equation}
The confinement scale $\lbar_p = \hbar/(m_p c)$ anchors the spatial extent (the pion Compton wavelength $\lbar_\pi = 1.41$~fm is anomalously large due to the Goldstone nature and does not set the internal structure scale).

\noindent Experimental: $r_\pi = 0.659 \pm 0.004$~fm. Error: 0.26\%, within 0.43$\sigma$.

%==================================================================
\section{Kaon Charge Radius}\label{sec:kaon}
%==================================================================

The kaon $K^+ = u\bar{s}$ involves quarks of different generations: $u$ at winding $w = 1$ and $\bar{s}$ at $w = 2$. The coherence must transfer \textit{between} Hopf fiber winding modes, governed by the fiber spectral gap.

\begin{proposition}[Kaon charge radius]\label{prop:kaon}
\begin{equation}\label{eq:kaon}
    r_K = \frac{1}{E_1(C_{10})} \times \lbar_p = \phig^2 \times \lbar_p = 0.5506~\text{fm}
\end{equation}
where $E_1(C_{10}) = 2 - 2\cos(\pi/5) = \phig^{-2}$ is the spectral gap of the Hopf decagonal fiber.
\end{proposition}

\noindent Experimental: $r_K = 0.560 \pm 0.031$~fm. Error: 1.7\%, within 0.30$\sigma$.

\begin{remark}[Scale-free prediction]
The ratio $r_\pi/r_K = \pi/\phig^2 = 1.200$ is a dimensionless prediction involving only structural constants. Experimental: $0.659/0.560 = 1.177 \pm 0.06$ (within $0.4\sigma$).
\end{remark}

%==================================================================
\section{Deuteron Charge Radius}\label{sec:deuteron}
%==================================================================

The deuteron is a proton--neutron bound state. Its charge radius decomposes as:
\begin{equation}\label{eq:deuteron}
    r_d^2 = r_p^2 + \langle r_n^2 \rangle + \frac{r_{\mathrm{str}}^2}{4}
\end{equation}
where $r_{\mathrm{str}}$ is the proton--neutron structural separation.

The strong nuclear exchange proceeds through pion (closure orbit phase mode) exchange between the two nucleons' shell supports. The factor $6 = 2 \times |V(P_3)|$: two nucleons, each with a 3-vertex support graph, exchanging pions through 3 shell-matched channels in each direction. Thus:
\begin{equation}
    r_{\mathrm{str}} = 6\,\pi\,\lbar_p
\end{equation}
and therefore:
\begin{equation}
    r_d = \sqrt{16 + 9\pi^2 - \tfrac{8}{3}} \times \lbar_p \approx 10.11 \times \lbar_p = 2.126~\text{fm}
\end{equation}

\noindent Experimental: $r_d = 2.12799 \pm 0.00074$~fm. Relative error: 0.10\%. The agreement is strong at the percent level, though outside experimental $1\sigma$ because the deuteron radius is measured very precisely. The residual discrepancy ($\sim\!0.002$~fm) is consistent in scale with an electromagnetic correction of order $\alpha\,\pi\,\lbar_p \approx 0.005$~fm, suggesting a natural direction for future extension once the interaction functional~\citep{SmartXIII} is derived directly from the closure functional rather than introduced at the two-body level.

%==================================================================
\section{The Electric Form Factor}\label{sec:formfactor}
%==================================================================

\subsection{Closed-form expression}

The resolvent at the boundary vertex of $P_3$ (Section~\ref{sec:resolvent}) gives the proton electric form factor as:
\begin{equation}\label{eq:formfactor}
    G_E(Q^2) = \frac{1 + 3z + z^2}{1 + 4z + 3z^2}
\end{equation}
where $z = (8/3)\,Q^2\,\lbar_p^2$. This is a rational function of $Q^2$: same functional class as the experimental dipole, monotonically decreasing, no spacelike zeros.

\subsection{Golden-ratio analytic structure}

The form factor factorises as:
\begin{equation}
    G_E = \frac{(z + \phig^{-2})(z + \phig^2)}{(z + 1)(z + 1/3)}
\end{equation}
The timelike zeros occur at the two algebraic scales $\phig^{-2}$ and $\phig^2$ --- the same golden-ratio pair that appears throughout the 600-cell shell and fiber spectral structure (the spectral gap of $P_5$, the Hopf fiber $C_{10}$, and the nearest-neighbour squared distance). The poles at $z = -1$ and $z = -1/3$ correspond to the nontrivial Laplacian eigenvalues of $P_3$.

\subsection{120-vertex resolvent}\label{sec:120vertex}

The full 120-vertex resolvent uses all 9 eigenvalue-multiplicity pairs of the 600-cell Laplacian:
\begin{equation}
    G_E^{(120)}(Q^2) = \frac{1}{120}\sum_k \frac{\mathrm{mult}_k}{1 + z\,\lambda_k}
\end{equation}
with $z = (2/9)\,Q^2\,\lbar_p^2$. This matches experiment within 0.4\% at $Q^2 = 0.1$~GeV$^2$ and asymptotes to $1/120 \approx 0.008$ (vs.\ $1/3$ for $P_3$), a 40-fold improvement at high $Q^2$.

%==================================================================
\section{The Unified Coherence Principle}\label{sec:unified}
%==================================================================

\begin{definition}[Spectral coherence radius]
For a particle with geometric home $G$, evaluation point $x$, and spectral channel $c$:
\begin{equation}\label{eq:unified}
    r = \xi_K(G, x, c) \times \frac{\hbar}{m_p\,c}
\end{equation}
where $\xi_K$ is the heat-kernel coherence length: the characteristic decay scale of the heat kernel $K(x,x;t) = \sum_k e^{-\lambda_k t}\,|\psi_k(x)|^2$ on $G$, measured in channel $c$.
\end{definition}

The resolvent (baryons) and diffusion length (mesons) are two moments of the same heat kernel: the resolvent is its Laplace transform, the diffusion length is its first moment. They are two functionals of the same object, applied to different geometric homes.

\begin{table}[ht]
\centering
\caption{Particle classes, geometric homes, and coherence mechanisms.}
\label{tab:unified}
\begin{tabular}{@{}llll@{}}
\toprule
Class & Geometric home & Coherence mechanism & Formula \\
\midrule
Charged baryon & Support graph $G(S)$ & Resolvent at boundary & $\xi_K = 12\,\mathrm{Tr}(L)/|V|$; $r = \sqrt{\xi_K}\,\lbar$ \\
Neutral baryon & Support graph $G(S)$ & Dipole eigenfunction & $\xi_K = -\mathrm{Var}(\deg)\,\Delta n^2$; $\langle r^2\rangle = \xi_K\,\lbar^2$ \\
Non-strange meson & Closure orbit $S^1$ & Phase half-period & $\xi_K = \pi$; $r = \xi_K\,\lbar_p$ \\
Strange meson & Hopf fiber $C_{10}$ & Inverse spectral gap & $\xi_K = \phig^2$; $r = \xi_K\,\lbar_p$ \\
\bottomrule
\end{tabular}
\end{table}

The selection of geometric home is not ad hoc: baryons \textit{are} closure basins (R1--R5), pions \textit{are} phase modes (Goldstone), kaons \textit{are} cross-fiber modes (strangeness). Each particle type selects its home by being what it is.

%==================================================================
\section{Sensitivity Analysis}\label{sec:sensitivity}
%==================================================================

\subsection{Sensitivity to the shell assignment}

Alternative assignments produce catastrophic failure:
\begin{itemize}
    \item $\{3,4,5\}$: $\mathrm{Tr}(L) = 4$ but $m_p/m_e = \phig^{48} \sim 10^{10}$, so $\lbar_p \sim 10^{-10}$~fm and $r_p \sim 10^{-10}$~fm.
    \item $\{1,2,3\}$: excluded by R2.
\end{itemize}
Only $\{2,3,4\}$ gives both the correct mass ratio \textit{and} the correct charge radius.

\subsection{Sensitivity to the operator}

\begin{table}[ht]
\centering
\caption{Alternative operators for the proton radius.}
\label{tab:operators}
\begin{tabular}{@{}llcc@{}}
\toprule
Operator & Value & $r_p$ (fm) & Error \\
\midrule
$|V|$ & 3 & 0.631 & 25\% \\
$|E|$ & 2 & 0.421 & 50\% \\
$\mathbf{Tr}(L)$ & $\mathbf{4}$ & $\mathbf{0.841}$ & $\mathbf{0.04\%}$ \\
$\max$ eigenvalue & 3 & 0.631 & 25\% \\
\bottomrule
\end{tabular}
\end{table}

For the proton support graph $P_3$, the coherence factor induced by the boundary resolvent reduces to the canonical invariant $\mathrm{Tr}(L) = 4$, giving $r = 4\,\lbar_p$.

%==================================================================
\section{Summary of Results}\label{sec:results}
%==================================================================

\begin{table}[ht]
\centering
\caption{All derived charge radii.}
\label{tab:results}
\begin{tabular}{@{}llcccc@{}}
\toprule
Observable & VFD formula & VFD value & Experiment & Error & $\sigma$ \\
\midrule
$r_p$ (charge) & $\sqrt{12\,\mathrm{Tr}(L)/|V|}\,\lbar_p$ & 0.8412~fm & $0.8409(4)$~fm & 0.04\% & 0.84 \\
$r_M$ (magnetic) & $\sqrt{12\,\mathrm{Tr}(L)/|V|}\,\lbar_p$ & 0.8412~fm & $0.851(26)$~fm & --- & 0.38 \\
$\langle r_n^2 \rangle$ & $-(8/3)\,\lbar_n^2$ & $-0.1176$~fm$^2$ & $-0.1161(22)$~fm$^2$ & 1.3\% & 0.69 \\
$r_\pi$ & $\pi \times \lbar_p$ & 0.661~fm & $0.659(4)$~fm & 0.26\% & 0.43 \\
$r_K$ & $\phig^2 \times \lbar_p$ & 0.551~fm & $0.560(31)$~fm & 1.7\% & 0.30 \\
$r_d$ & $\sqrt{16{+}9\pi^2{-}8/3}\,\lbar_p$ & 2.126~fm & $2.12799(74)$~fm & 0.10\% & 3.0 \\
\bottomrule
\end{tabular}
\end{table}

\noindent Six observables derived within the framework from one geometric principle, with zero fitted parameters. Five lie within $1\sigma$ of experiment. The deuteron shows strong relative agreement (0.10\%) but lies outside $1\sigma$ because the experimental uncertainty is very small.

%==================================================================
\section{Discussion}\label{sec:discussion}
%==================================================================

\subsection{What is established}

\begin{itemize}
    \item The proton charge radius is derived from the boundary resolvent on the uniquely determined support graph $P_3$, for which the resulting coherence factor reduces to the graph invariant $\mathrm{Tr}(L(P_3)) = 4$.
    \item The neutron charge radius uses the dipole eigenfunction of the same graph, with amplitude $\mathrm{Var}(\deg) = 2/9$.
    \item The pion and kaon use different coherence channels: phase ($\pi$) and spectral gap ($\phig^2$), both derived from the 600-cell geometry.
    \item The form factor is a closed-form rational function with golden-ratio zeros.
    \item All six results are unified by one principle: heat-kernel coherence length.
\end{itemize}

\subsection{What is not established}

\begin{itemize}
    \item A first-principles derivation of \textit{why} the closure boundary determines the electromagnetic response. Boundary dominance is proved from the barrier structure (Section~\ref{sec:boundary_dominance}) but the connection from the closure functional to the electromagnetic current is a modelling step.
    \item The full form factor curve at $Q^2 > 0.5$~GeV$^2$. The 120-vertex resolvent improves over $P_3$ but still drops too slowly compared to the experimental dipole.
    \item Extension to heavy mesons ($D$, $B$), whose charge radii would test the cross-generation formula at $\Delta w = 2$ (predicted: $r_B \approx [1/(3-\phig)] \times \lbar_p$).
    \item Loop corrections to the charge radii, which require the closure residual $\delta U_{\mathrm{rel}}$ (Papers~XIX--XX).
\end{itemize}

\subsection{Relation to existing framework results}

The charge radii use the same graph invariants that enter the mass formula~\citep{SmartIV}: $\mathrm{Tr}(L)/|V|$ is the first cumulant coefficient (first-order mass correction), $\mathrm{Var}(\deg)$ is the second cumulant (second-order correction). The mass formula and radius formula are complementary projections of the same heat-kernel expansion: the mass uses the per-vertex average $\mathrm{Tr}(L)/|V|$, the radius uses the total $\mathrm{Tr}(L)$.

%==================================================================
\section{Reproducibility}\label{sec:reproducibility}
%==================================================================

Every claim reduces to a computation verifiable with \texttt{numpy}:
\begin{enumerate}
    \item \textbf{Proton radius.} The factor 4 is $\mathrm{Tr}(L(P_3)) = 0 + 1 + 3$, verifiable by diagonalising a $3 \times 3$ integer matrix.
    \item \textbf{Neutron radius.} $\mathrm{Var}(\deg) = 2/9$ from degree sequence $[1,2,1]$; $-(8/3) \times (0.21031)^2 = -0.1179$.
    \item \textbf{Form factor.} The rational function $(1+3z+z^2)/(1+4z+3z^2)$ with zeros at $(-3\pm\sqrt{5})/2$.
    \item \textbf{600-cell spectrum.} The $120 \times 120$ Laplacian is constructible from standard vertex coordinates.
\end{enumerate}

\noindent The accompanying repository (\url{https://github.com/vfd-org/vfd-crystallisation}) includes \texttt{build\_600cell.py} (600-cell construction and spectral verification), \texttt{verify\_hadron\_radii.py} (automated verification of all numerical claims), and \texttt{600cell\_data.npz} (precomputed vertex and spectral data).

%==================================================================
\section{Conclusion}\label{sec:conclusion}
%==================================================================

Six hadron charge radii emerge from the spectral geometry of the 600-cell, unified by one principle: the charge radius is the heat-kernel coherence length of the electromagnetic response, computed on the geometric sub-object where the particle lives. The proton radius $r_p = \sqrt{12\,\mathrm{Tr}(L)/|V|}\,\hbar/(m_p c) = 4\,\hbar/(m_p c) = 0.8412$~fm (0.04\% error) follows from the boundary resolvent on the uniquely determined support graph $P_3$. The kaon radius $r_K = \phig^2 \times \hbar/(m_p c)$ connects the charge radius to the Hopf fiber spectral gap and the generation structure. The form factor is a rational function with golden-ratio zeros, matching experiment within 0.4\% at low $Q^2$.

The results add six independent observables to the VFD closure programme, which previously derived particle masses~\citep{SmartIV, SmartV}, coupling constants~\citep{SmartXXII}, and gauge structure~\citep{SmartVI, SmartXXII} from the same 600-cell geometry. The charge radii use the same graph invariants as the mass formula, confirming internal consistency. The framework provides a unified geometric account of masses, coupling constants, and spatial structure within the closure programme from one geometric object and zero fitted parameters.

%==================================================================
\bibliographystyle{plainnat}
\bibliography{references}

\end{document}
